\paragraph{Data.} All inputs are public administrative aggregates (daily
counts of service requests by category) and public weather observations.
No individual-level, household-level, or address-level information is
acquired, stored, or processed; the acquisition layer requests
server-side aggregation precisely so that record-level attributes
(including free-text descriptions and locations) never enter the
pipeline.

\paragraph{People inside the data.} Each request count summarizes
residents' attempts to obtain public services. Two ethical asymmetries
follow. First, reported volume under-represents communities with lower
propensity or ability to report \citep{kontokosta2017equity,
kontokosta2021bias}; optimizing service against reported volume can
therefore divert resources from quiet need toward vocal demand. This
study's city-level action space cannot reallocate across neighbourhoods,
which contains but does not eliminate the concern: the family-level
analogue (e.g., systematically deprioritizing housing complaints relative
to noise) is measured and reported per policy. Second, the priority
weights are normative inputs; we publish them as configuration, sweep
them, and explicitly assign their real-world selection to accountable
human decision-makers.

\paragraph{Intended and unintended use.} The system evaluated here is a
research benchmark. It is not a deployable dispatch tool, has no city
involvement or endorsement, and the manuscript states this in every
section where simulation results appear. Misuse risk (treating simulated
capacity numbers as operational guidance) is mitigated by classifying all
operational parameters as sensitivity-only and reporting ranges rather
than recommendations.

\paragraph{Human oversight.} The intended human-AI division of labour ---
forecasts and allocation suggestions with uncertainty made visible,
final allocation authority retained by planners --- is described in the
discussion and is a design commitment of the broader research programme,
not an evaluated capability of this study.
